\section*{Capítulo \ref{ch:g12:continental-crust} --- A crosta continental}

\begin{solution}{exo:g12:continental-crust:1}
\omterm{def:g12:continental-crust:continental-crust}{Crosta continental}: de 30 a \qty{70}{km}, cerca de \qty{41}{km} em média;
rochas graníticas, gnaisses e outras rochas metamórficas e \omterm{def:g11:igneous-rocks:volcanic-plutonic}{plutônicas} sob
uma cobertura sedimentar; cerca de \qty{61}{\percent} de sílica em média;
cerca de \qty{2830}{kg/m^3}; até 4.03 bilhões de anos (grãos de 4.4).
\omterm{def:g11:ocean-ridges:oceanic-lithosphere}{Crosta oceânica}: cerca de \qty{7}{km}; basalto e \omterm{def:g11:ocean-ridges:gabbro}{gabro}; cerca de
\qty{50}{\percent} de sílica; cerca de \qty{3000}{kg/m^3}; menos de 200
milhões de anos.
\end{solution}

\begin{solution}{exo:g12:continental-crust:2}
\qty{60.6}{\percent}: \omterm{def:g11:igneous-rocks:felsic-mafic}{intermediária}, como um \omterm{def:g11:igneous-rocks:rock-names}{diorito}. \qty{66.6}{\percent}:
\omterm{def:g11:igneous-rocks:felsic-mafic}{félsica}, como um granito.
\end{solution}

\begin{solution}{exo:g12:continental-crust:3}
A \omterm{def:g11:ocean-ridges:oceanic-lithosphere}{crosta oceânica}, com cerca de \qty{3000}{kg/m^3}, assenta sobre o manto
litosférico, e a \omterm{def:g11:ocean-ridges:oceanic-lithosphere}{litosfera oceânica} velha e fria, como um todo, é densa o
bastante para afundar. A \omterm{def:g12:continental-crust:continental-crust}{crosta continental}, com cerca de
\qty{2830}{kg/m^3}, muito mais leve que o manto (cerca de
\qty{3380}{kg/m^3}), flutua demais: um continente que chega a uma fossa
interrompe a \omterm{def:g11:subduction-zones:subduction}{subducção}.
\end{solution}

\begin{solution}{exo:g12:continental-crust:4}
(1) Um núcleo antigo e estável de um continente, que não se deforma mais.
(2) Mais de um bilhão de anos. (3) Rochas metamórficas e \omterm{def:g11:igneous-rocks:volcanic-plutonic}{plutônicas}, as
\omterm{def:g12:continental-crust:root}{raízes} profundas de antigas montanhas.
\end{solution}

\begin{solution}{exo:g12:continental-crust:5}
$4030/190 \approx 21$ vezes. O fundo oceânico é produzido nas dorsais e
destruído nas \omterm{def:g11:subduction-zones:subduction}{zonas de subducção}: nada dele sobrevive muito além de 200
milhões de anos.
\end{solution}

\begin{solution}{exo:g12:continental-crust:6}
(1) $0.306 \times 210.4 \approx \num{64.4}\text{ milhões de }\unit{km^2}$.
(2) $210.4 - 64.4 = \num{146.0}\text{ milhões de }\unit{km^2}$.
(3) $4\pi \times 6371^2 \approx \qty{5.10e8}{km^2}$; $210.4/510.1 \approx
0.41$: \qty{41}{\percent}.
\end{solution}

\begin{solution}{exo:g12:continental-crust:7}
(1) $41100 \times 2830 \approx \qty{1.16e8}{kg}$. (2) $7100 \times 3000 =
\qty{2.13e7}{kg}$. (3) Cerca de 5.5 vezes.
\end{solution}

\begin{solution}{exo:g12:continental-crust:8}
\omterm{def:g12:continental-crust:root}{Raiz}: $2.5 \times 2830/550 \approx \qty{12.9}{km}$. A crosta é mais espessa
de $2.5 \times 3380/550 \approx \qty{15.4}{km}$: cerca de \qty{50}{km} ao
todo, dentro dos 50 a \qty{70}{km} encontrados sob os cinturões de
montanhas.
\end{solution}

\begin{solution}{exo:g12:continental-crust:9}
(1) Cerca de $0.103 + 0.100 + 0.102 \approx 0.30$: \qty{30}{\percent}.
(2) A faixa mais jovem, de 0 a 10 milhões de anos (cerca de
\qty{10}{\percent}), quase empatada com a de 20 a 30. (3) O fundo antigo foi
destruído nas \omterm{def:g11:subduction-zones:subduction}{zonas de subducção}; quanto mais antigo o fundo, mais dele
desapareceu.
\end{solution}

\begin{solution}{exo:g12:continental-crust:10}
$120 = 0.694\,h - 0.306 \times 100$, logo $h = 150.6/0.694 \approx
\qty{217}{m}$. Contar só as terras emersas deixa de fora as bordas baixas e
submersas dos continentes, e a média sai muito maior.
\end{solution}

\begin{solution}{exo:g12:continental-crust:11}
A ilha: um \omterm{def:g11:ocean-ridges:gabbro}{gabro} sob lavas e um \omterm{def:g11:seismic-structure:moho}{Moho} a \qty{8}{km}, \omterm{def:g11:ocean-ridges:oceanic-lithosphere}{crosta oceânica} (a
altitude da ilha não muda isso). A planície: gnaisse e um \omterm{def:g11:seismic-structure:moho}{Moho} a
\qty{38}{km}, \omterm{def:g12:continental-crust:continental-crust}{crosta continental}.
\end{solution}

\begin{solution}{exo:g12:continental-crust:12}
(1) Exato: $\tfrac{4}{3}\pi(6371^3 - 6329.9^3) \approx
\qty{2.0829e10}{km^3}$; aproximado: $4\pi \times 6371^2 \times 41.1
\approx \qty{2.0964e10}{km^3}$. (2) A aproximação é grande demais em cerca
de \qty{0.65}{\percent}, perto de $H/R \approx 0.0065$: para uma camada
fina, o erro é desprezível.
\end{solution}

\begin{solution}{exo:g12:continental-crust:13}
(1) $60.6 = (66.6 + 2x)/3$, logo $x = (181.8 - 66.6)/2 = \qty{57.6}{\percent}$.
(2) \omterm{def:g11:igneous-rocks:felsic-mafic}{Intermediária}. (3) A \omterm{def:g12:continental-crust:continental-crust}{crosta continental} profunda é menos rica em sílica
que a crosta superior: mais minerais \omterm{def:g11:igneous-rocks:felsic-mafic}{máficos}, mais rica em ferro, magnésio e
cálcio.
\end{solution}

\begin{solution}{exo:g12:continental-crust:14}
(1) $3 \times 10^8/3 = 10^8$ anos, 100 milhões de anos. (2) A mesma ordem de
grandeza que as idades do fundo oceânico: três quartos com menos de 100
milhões de anos, nada com mais de cerca de 190. (3) Ele deve ser destruído
tão depressa quanto é produzido: volta a afundar no manto nas \omterm{def:g11:subduction-zones:subduction}{zonas de
subducção}.
\end{solution}

\begin{solution}{exo:g12:continental-crust:15}
(1) $8.6 \times 10^9 / 4.0 \times 10^9 \approx \qty{2.2}{km^3}$ por ano.
(2) Uma taxa constante é só uma possibilidade: a crosta pode ter crescido
depressa no início e sido reciclada desde então, ou ter crescido tarde. As
rochas e os grãos mais antigos provam que alguma crosta existia cedo; eles
não podem dizer quanto, porque a crosta antiga pode ser erodida ou levada de
volta ao manto e deixar poucos vestígios.
\end{solution}

\begin{solution}{pb:g12:continental-crust:1}
\textbf{1.} $4\pi \times 6371^2 \approx \qty{5.10e8}{km^2}$.

\textbf{2.} $210.4/510.1 \approx 0.41$: \qty{41}{\percent}.

\textbf{3.} $0.306 \times 210.4 \approx \num{64.4}\text{ milhões de }\unit{km^2}$.

\textbf{4.} $510.1 - 210.4 \approx \num{299.7}\text{ milhões de }\unit{km^2}$,
\qty{59}{\percent}.

\textbf{5.} As bordas dos continentes ficam sob mares rasos, nas
plataformas continentais.

\textbf{6.} A espessura, \qty{41}{km}, é muito pequena em comparação com o
raio, \qty{6371}{km}: a curvatura muda o volume em menos de um por cento.

\textbf{7.} $210.4 \times 10^6 \times 41.1 \approx \qty{8.65e9}{km^3}$.

\textbf{8.} \qty{8.65e18}{m^3}.

\textbf{9.} $210.4 \times 10^6 \times 36 \approx \qty{7.57e9}{km^3}$.

\textbf{10.} $(8.65 - 7.57)/7.57 \approx 0.14$: \qty{14}{\percent}.

\textbf{11.} $8.65 \times 10^{18} \times 2830 \approx \qty{2.45e22}{kg}$.

\textbf{12.} $41100 \times 2830 \approx \qty{1.16e8}{kg}$.

\textbf{13.} $299.7 \times 10^6 \times 7.1 \approx \qty{2.13e9}{km^3}$, isto
é, \qty{2.13e18}{m^3}; vezes 3000: cerca de \qty{6.4e21}{kg}.

\textbf{14.} A \omterm{def:g12:continental-crust:continental-crust}{crosta continental}, cerca de 3.8 vezes mais pesada: ela é
quase seis vezes mais espessa, o que compensa a sua área e a sua densidade
menores.

\textbf{15.} Ela é feita de rochas graníticas, ricas em quartzo e em
feldspatos, os minerais formadores de rochas mais leves; a \omterm{def:g11:ocean-ridges:oceanic-lithosphere}{crosta oceânica} é
feita de basalto e de \omterm{def:g11:ocean-ridges:gabbro}{gabro}, ricos em piroxênio e em olivina.

\textbf{16.} $2.45 \times 10^{22}/5.97 \times 10^{24} \approx 0.0041$.

\textbf{17.} $3 \times 0.41 \times (41.1/6371) \times (2830/5513) \approx
0.0041$: o mesmo valor.

\textbf{18.} $6.4 \times 10^{21}/5.97 \times 10^{24} \approx 0.0011$: cerca
de um décimo de por cento.

\textbf{19.} Cerca de \qty{0.4}{\percent}.

\textbf{20.} A \omterm{def:g12:continental-crust:continental-crust}{crosta continental} é cerca de \textbf{\qty{0.4}{\percent} da
massa da Terra}.
\end{solution}
